% SHARD-ID: RC-10-OPEN-DEAD
% SHARD-TITLE: Open Problems and Dead Routes
% SHARD-SUMMARY: Lists the open steps in order, each as a tagged conjecture, and the routes known not to work.
% SHARD-SUMMARY: Mirrors the next-steps list of HANDOFF.md and the dead-route discussion of the worklogs.
% SHARD-KEYWORDS: open problems, dead routes, next steps, Stinespring, assemble over p, Ramanujan for Kraus
% SHARD-NOTES: none
% SHARD-SCRIPTS: none

\section{Open Problems and Dead Routes}
\label{sec:open}

Dead routes are results. Both lists mirror \texttt{HANDOFF.md}, which is
the live version.

\subsection{Open problems, in order}

\begin{conjecture}[Ramanujan without unitarity]
\label{conj:kraus-ramanujan}
\claimstatus{conj:kraus-ramanujan}{open}
There is a Ramanujan-type reading of the poles of $\zetaT$ for non-unitary
Kraus families, replacing the relation $\edgeev\edgeev' = \qs$ of the unitary case,
and canonical form $\sum_k \Kr{k}^\dagger \Kr{k} = 1$ constrains the poles.
\end{conjecture}

What is known: \cref{thm:qihara-general} replaces $\qs\uz^2$ by the
operator $\Dop$ and $\uz\,\Sig$ by the deformed channel $\Aop$, both
built from $(1-\uz^2\Ad(\Kr{k}^\dagger \Kr{k}))^{-1}$, so no quadratic relation pairs
the eigenvalues of $\Hash$ with those of $\Sig$. Cheap next step: a random-MPS
pole-radius hunt (bond dimension $2$ to $4$, physical dimension $2$) to see
whether the poles cluster on one circle at all.

\paragraph{Stinespring form of the compressed semigroup.} \Cref{conj:stinespring-jump-form}
(stated in \cref{sec:riemann-channel}): a Stinespring form of $\Zt$ on $\KS$ whose Kraus
(jump) operators are the prime dilations.

The single-Kraus lift exists trivially. A Holevo-form covariant Lindbladian
with the von Mangoldt jump measure lives on the full regular representation,
has symbol $\zeta^t$, and is not the compressed object. A negative answer with
a reason is a result. Suggested route (HANDOFF step 2): write $\Zt$ explicitly
through Cauchy kernels and test a generator of Holevo jump form with jumps at
$k\log \pp$ on $\KS$, blind, in a second model family.

\paragraph{Hecke identification of the Weil--LPS spectra.} \Cref{conj:weil-lps-hecke}
(stated in \cref{sec:weil-lps}): the joint spectra of the commuting channels
$\PhiWL{\qs}^{\pm}$ are Hecke eigenvalues of weight-$2$ forms for the quaternion algebra
ramified at $\{2,\infty\}$, via Jacquet--Langlands.

Needs LMFDB and the Jacquet--Langlands bookkeeping; the visible structure
(joint eigenvalue pairs, small-degree algebraic integers, multiplicities equal
to those of irreducibles of $\PSL(\Fp)$ in $\Wpm\otimes\overline{\Wpm}$)
is consistent with it.

\begin{conjecture}[Assembling over $\pp$]
\label{conj:assemble-over-p}
\claimstatus{conj:assemble-over-p}{open}
The Weil--LPS channels for different $\pp$ assemble into one object whose rings
are the primes and whose transfer spectrum is the zeros of $\zeta$.
\end{conjecture}

This is the parent campaign's DG-GLOBAL question in this guise. Nothing is
known beyond the obstruction recorded below.

\begin{conjecture}[The Phantasm has both halves]
\label{conj:phantasm-both-halves}
\claimstatus{conj:phantasm-both-halves}{open}
There is a $\mathbb Z_2$-graded Lindbladian (or a type~III generalisation at
$\invtemp=1$) whose stationary state is the Bost--Connes KMS state, with entanglement
Hamiltonian $\log N$, and whose odd relaxation modes are the nontrivial zeros: the object
having both halves of \cref{obs:complementary-halves}. The rebound-state equation of the
SHW renewal picture with a \emph{mixed} rebound state is its finite-dimensional shadow;
the vacuum rebound state is excluded by \cref{prop:vacuum-renewal-trivial-entanglement}.
\end{conjecture}

This replaces item 0b of the reframe programme. The fixed point constrains the Lindbladian
only partially, and at $\invtemp=1$ the fixed point is not normal
(\cref{prop:normal-kms-gibbs}), so $\KS$ with its pure vacuum cannot be the whole bond; the
identification of $\KS$ with the Bost--Connes bond that \cref{sec:riemann-channel} records
as missing is the first obstacle. A second open item from the same campaign is the
unrestricted infinite parity question, \cref{obs:parity-infinite-open}.

\begin{conjecture}[Galois-graded bond at finite level]
\label{conj:galois-graded-bond}
\claimstatus{conj:galois-graded-bond}{open}
The model space of a congruence subgroup $\Gamma_0(N)$, whose scattering determinant
involves the Dirichlet $L$-functions mod $N$, carries the zeros of $L(s,\chi)$ organised
by the characters $\chi$ of $(\mathbb Z/N)^\times$, with the Shor map $\shor{a}$ as the
level-$N$ Galois symmetry acting on the bond; the graded bond of
\cref{obs:character-sector-grading} is realised at finite level.
\end{conjecture}

This is a computation, not yet done; it is the Riemann-side test of the
trivial-versus-nontrivial-character reading of the grading that the Artin--Schreier
super-transfer matrix exhibits.

\subsection{Dead routes}

\begin{observation}[Prime-by-prime ans\"atze are blind to zeros]
\label{obs:prime-by-prime-blind}
\claimstatus{obs:prime-by-prime-blind}{sketched}
The direct sum over primes of rotations on circles of circumference $\log \pp$
has trace $\sum_n\vM(n)\,\delta(s-\log n)$ by Poisson summation and
dynamical zeta equal to the Euler product; for any finite set $S$ of primes
its spectrum is the set of poles of the partial Euler product. The same holds
for product states and commuting dilations. Zeros exist only in the infinite
product, so no prime-by-prime ansatz can see them.
\end{observation}

Recorded 2026-09-12 from the circle-flow discussion. Two remarks kept: in the
Witt convention $L_1 - L_{-1}$ is the dilation flow (source and sink) and
$L_0$ the rotation; the only relation between primes is additive.

\begin{observation}[Symmetry protection fixes signs, not moduli]
\label{obs:spt-protection-is-sign-data}
\claimstatus{obs:spt-protection-is-sign-data}{sketched}
The content of an L-function that has the three properties of a
symmetry-protected topological invariant (a phase, quantised, robust under
symmetric perturbation) is its sign data: the root number as a product of
local $\epsilon$-factors with the product formula as anomaly cancellation,
the sign of a quadratic Gauss sum (the Arf invariant over $\mathbb F_2$, the
Gauss--Milgram central-charge phase over $\mathbb Z/p$), and Hilbert-symbol
reciprocity as a local-to-global cocycle condition. The Riemann hypothesis is
modulus data. Moduli are not protected: a Ramanujan graph ceases to be
Ramanujan under a small symmetric perturbation, and the cubic Cayley graphs
of $\PSL(\Fp)$ on the generators of $\PSL(2,\mathbb Z)$ fail the Ramanujan
bound for $\pp\ge 11$ while keeping every symmetry. Cohomological protection
therefore cannot deliver RH.
\end{observation}

Recorded 2026-09-12 from the SPT side quest. What survives of the
bulk--boundary lesson is the positivity form of Weil's criterion: the
reflection $\zero\mapsto1-\bar\zero$ pairs the mode at a zero with the mode at
its mirror image, RH says every mode is its own partner, and Weil's form is
the pairing of each mode with its partner, written through the explicit
formula as a form on the primes. Its finite-dimensional, arbitrary-Kraus
version is the subject of \cref{sec:weil-positivity}. Two relabellings, not new
content: the Weil--LPS channels live on the edge Hilbert space of the
$\mathbb Z/p\times\mathbb Z/p$ cluster-state phase (the Heisenberg group is
that edge's projective representation and $\SL(\Fp)$ its automorphism group,
acting by the Weil representation); and Bost--Connes at $\invtemp=1$ is a
critical point of a spontaneous breaking of the Galois symmetry, not a
symmetry-protected phase.

Further routes closed or flagged in the founding session: the functional
integral framing of Selberg as a continuum limit (TJO: too far); calling
$-\log \Tmat$ the entanglement Hamiltonian (an error, corrected in
\cref{obs:hilbert-polya-hermitian}); an algebraic-integer recognition of the
Weil--LPS joint eigenvalues from five-digit output (tautological, since the
adjacency matrix is integral); the Dwork transfer operator as a route to the
Riemann hypothesis for curves (\cref{obs:as-nonquadratic}).
